% Generated draft: requires mathematical and source review.
\documentclass[11pt]{article}
\usepackage[margin=25mm]{geometry}
\usepackage{fontspec}
\setmainfont{Latin Modern Roman}
\usepackage{amsmath,amssymb,mathtools,unicode-math}
\setmathfont{Latin Modern Math}
\usepackage{enumitem,needspace,xurl,hyperref}
\setlength{\parindent}{0pt}
\setlength{\parskip}{5pt}
\setlength{\emergencystretch}{3em}
\newcommand{\N}{\mathbb{N}}
\newcommand{\Z}{\mathbb{Z}}
\newcommand{\Q}{\mathbb{Q}}
\newcommand{\R}{\mathbb{R}}
\newcommand{\C}{\mathbb{C}}
\newcommand{\F}{\mathbb{F}}
\newcommand{\A}{\mathbb{A}}
\newcommand{\PP}{\mathbb{P}}
\newcommand{\E}{\mathbb{E}}
\newcommand{\eps}{\varepsilon}
\newcommand{\dd}{\,\mathrm{d}}
\begin{document}
\section*{UnsolvedMath ID: 20000251}
\textbf{Source problem code:} catalogue entry
\subsection{Definitions and mathematical statement}
The supplied reconstruction of Madhav Nori's Question 14 asks: if $X$ is a smooth affine variety of dimension $d$ and $\alpha\in F^{d-1}K_0(X)$, must there be a vector bundle $E$ of rank $d-1$ with $[E]=\alpha$? Here $K_0$ is the Grothendieck group of vector bundles; under the topological filtration convention, $F^iK_0$ is generated by coherent classes supported in codimension at least $i$, using smoothness to identify coherent and vector-bundle $K$-theory. This filtration convention is not defined in the damaged excerpt.

For $d\geq2$, every such $\alpha$ has rank zero, while $[E]$ has rank $d-1>0$. Thus the printed equality is rank-inconsistent, as the report emphasizes. Its proposed corrected question is explicitly different:
\[
 [E]-(d-1)[\mathcal O_X]=\alpha,\qquad\operatorname{rank}E=d-1.
\]
No assertion is made that this reconstruction was the intended printed statement.

The report says that a rank-$d$ representative $W$ of $d[\mathcal O_X]+\alpha$ always exists by stable realization on a smooth affine $d$-fold. It identifies corrected solvability with existence of some such representative splitting a free line:
\[
 \exists E\ \Longleftrightarrow\
 \exists W\text{ of rank }d:\ [W]=d[\mathcal O_X]+\alpha,
 \quad W\cong\mathcal O_X\oplus E.
\]
This existential formulation does not assume cancellation or that an arbitrary chosen representative splits. Over an algebraically closed field, the report attributes to Serre stable realization and the Murthy--Krishna top-rank splitting theorem the criterion
\[
 \text{corrected solvability}\quad\Longleftrightarrow\quad
 c_d(\alpha)=0\quad\text{in }\operatorname{CH}^d(X).
\]
The precise dimension and characteristic hypotheses of those cited results are not supplied.

For a surface and a rank-zero class, the report defines the determinant defect
\[
 \delta_2(\alpha)=\alpha-\bigl([\det(\alpha)]-[\mathcal O_X]\bigr).
\]
It vanishes exactly when the corrected rank-one representation is given by the determinant line bundle. The report refutes the corrected universal assertion over $\C$ using the smooth affine surface
\[
 X=\{(x,y,z)\in\A^3_{\C}:1+x^4+y^4+z^4=0\}.
\]
It claims $\operatorname{CH}_0(X)\ne0$; for a nonzero zero-cycle class $\zeta$, the cycle class $\alpha=\operatorname{cyc}_K(\zeta)\in F^2K_0(X)$ satisfies $c_2(\alpha)=-\zeta\ne0$, and therefore is not $[L]-[\mathcal O_X]$ for any line bundle $L$. No explicit representative cycle or proof of the nonzero Chow group is supplied. These claims address $d\geq2$; the original excerpt's unspecified low-dimensional range needs separate interpretation.
\subsection{Short English statement}
The printed bundle-class equality has incompatible ranks in dimension at least two. The report separately studies a rank-corrected version, reducing it to splitting a free line from some rank-$d$ representative and claiming a top-Chern criterion over algebraically closed fields. A supplied affine-surface example reportedly disproves even that corrected universal statement.
\subsection{Sources}
\begin{itemize}
\item Local registry: \url{data/unsolvedmath/problems.json}, permanent ID $20000251$, source code \texttt{catalogue entry}; damaged original text, supplied reconstruction and complete research summary. Statement status reconstructed and unverified.
\item AIM workshop ``Projective modules and $A^1$-homotopy theory'', Question 14, attributed to Madhav Nori: \url{https://aimath.org/pastworkshops/projectiveA1problems.pdf}. Canonical location \texttt{aim-algebraic-geometry-notes.json}, \texttt{notes[250]}; original literature field empty. No PDF was opened here.
\item Supplied attributions: Serre stable realization and the Murthy--Krishna top-rank splitting theorem; no full bibliographic references, theorem statements or characteristic hypotheses supplied.
\item Supplied machine-generated research attempt 1: counterexample; status refuted at the run, registry solved and classification \texttt{SOLVED-IN-LITERATURE}; rigor metadata ``proved'', novelty confidence low. The non-peer-reviewed report requires independent verification before citation.
\end{itemize}
\end{document}
